\documentclass[11pt, a4paper]{article}\usepackage[]{graphicx}\usepackage[]{xcolor}
% maxwidth is the original width if it is less than linewidth
% otherwise use linewidth (to make sure the graphics do not exceed the margin)
\makeatletter
\def\maxwidth{ %
  \ifdim\Gin@nat@width>\linewidth
    \linewidth
  \else
    \Gin@nat@width
  \fi
}
\makeatother

\definecolor{fgcolor}{rgb}{0.345, 0.345, 0.345}
\newcommand{\hlnum}[1]{\textcolor[rgb]{0.686,0.059,0.569}{#1}}%
\newcommand{\hlsng}[1]{\textcolor[rgb]{0.192,0.494,0.8}{#1}}%
\newcommand{\hlcom}[1]{\textcolor[rgb]{0.678,0.584,0.686}{\textit{#1}}}%
\newcommand{\hlopt}[1]{\textcolor[rgb]{0,0,0}{#1}}%
\newcommand{\hldef}[1]{\textcolor[rgb]{0.345,0.345,0.345}{#1}}%
\newcommand{\hlkwa}[1]{\textcolor[rgb]{0.161,0.373,0.58}{\textbf{#1}}}%
\newcommand{\hlkwb}[1]{\textcolor[rgb]{0.69,0.353,0.396}{#1}}%
\newcommand{\hlkwc}[1]{\textcolor[rgb]{0.333,0.667,0.333}{#1}}%
\newcommand{\hlkwd}[1]{\textcolor[rgb]{0.737,0.353,0.396}{\textbf{#1}}}%
\let\hlipl\hlkwb

\usepackage{framed}
\makeatletter
\newenvironment{kframe}{%
 \def\at@end@of@kframe{}%
 \ifinner\ifhmode%
  \def\at@end@of@kframe{\end{minipage}}%
  \begin{minipage}{\columnwidth}%
 \fi\fi%
 \def\FrameCommand##1{\hskip\@totalleftmargin \hskip-\fboxsep
 \colorbox{shadecolor}{##1}\hskip-\fboxsep
     % There is no \\@totalrightmargin, so:
     \hskip-\linewidth \hskip-\@totalleftmargin \hskip\columnwidth}%
 \MakeFramed {\advance\hsize-\width
   \@totalleftmargin\z@ \linewidth\hsize
   \@setminipage}}%
 {\par\unskip\endMakeFramed%
 \at@end@of@kframe}
\makeatother

\definecolor{shadecolor}{rgb}{.97, .97, .97}
\definecolor{messagecolor}{rgb}{0, 0, 0}
\definecolor{warningcolor}{rgb}{1, 0, 1}
\definecolor{errorcolor}{rgb}{1, 0, 0}
\newenvironment{knitrout}{}{} % an empty environment to be redefined in TeX

\usepackage{alltt}

\usepackage[top = 0.6 in, bottom = 0.6 in, left = 0.8 in, right = 0.8 in]{geometry}
\usepackage{amsmath, amssymb, amsfonts}

\allowdisplaybreaks[4]

\usepackage{enumerate}
\usepackage{array}
\usepackage{multirow}
\usepackage{dingbat}
\usepackage{fontawesome5}
\usepackage{tasks}
\usepackage{bbding}
\usepackage{twemojis}
% how to use bull's eye ----- \scalebox{2.0}{\twemoji{bullseye}}
\usepackage{fontspec}
\usepackage{fancyhdr}
\usepackage{lastpage}

\pagestyle{fancy}
\fancyhf{}
\fancyfoot[C]{Page \thepage\ of \pageref{LastPage}}
\renewcommand{\headrulewidth}{0pt}

\title{MSMS 408 : Practical 11}
\author{{\fontsize{22}{20} \benyorithafont Ananda Biswas} \\[1em] Exam Roll No. : 24419STC053}
\date{\today}

\newfontface\benyorithafont{Benyoritha-G334D.otf}

\newfontface\myfont{Myfont1-Regular.ttf}[LetterSpace=0.05em]

\newfontface\cbfont{CaveatBrush-Regular.ttf}
% how to use --- \myfont --write text here--
\IfFileExists{upquote.sty}{\usepackage{upquote}}{}
\begin{document}

\maketitle


\section*{\faArrowAltCircleRight[regular] \textcolor{blue}{Question}}

\hspace{1cm} Consider the data-set representing repair times of machines (in hours) $$X = \{2.1, 1.5, 3.2, 0.8, 1.9, 2.7, 1.1, 0.9\}.$$ Assume that the data come from an Exponential distribution with parameter $\theta$.

\begin{enumerate}[(i)]
\item Use the sample to provide a maximum likelihood estimate of $\theta$.

\item Employ the parametric bootstrap method to generate $B = 1000$ bootstrap samples from the exponential distribution with the estimated parameter plugged in and compute the bootstrap estimates $\widehat{\theta^*}$.

\item Using these bootstrap estimates, estimates the bias, standard error and variance of $\widehat{\theta}$.

\item Construct $95\%$ confidence interval for $\theta$.

\item Comment on the shape of the bootstrap distribution.
\end{enumerate}

\section*{\faArrowAltCircleRight[regular] \textcolor{blue}{Theory}}

The maximum likelihood estimate of $\theta$ is $$\widehat{\theta_{\text{MLE}}} = \dfrac{1}{\bar{x}}.$$

For bootstrapping, we proceed as follows :

\begin{enumerate}[Step I.]
\item Suppose we have $x_1, x_2, \ldots, x_n \sim f(x, \theta)$ where functional form of $f$ is known but $\theta$ is unknown.

\item Obtain an estimate of $\theta$ as $\hat{\theta} = \dfrac{1}{\bar{x}}$.

\item Plug in $\hat{\theta}$ in place of $\theta$ and generate a sample of size $n$ as $x_1^*, x_2^*, \ldots, x_n^* \sim f(x, \hat{\theta})$.

\item Based on bootstrapped sample $\{x_1^*, x_2^*, \ldots, x_n^*\}$, calculate $\widehat{\theta_{j}^{*}}$.

\item Repeat steps III and IV both $B$ times to get $$\widehat{\theta_{1}^{*}}, \widehat{\theta_{2}^{*}}, \ldots, \widehat{\theta_{B}^{*}}.$$
\end{enumerate}

$\{\widehat{\theta_{1}^{*}}, \widehat{\theta_{2}^{*}}, \ldots, \widehat{\theta_{B}^{*}}\}$ can be thought as sample from the sampling distribution of $\hat{\theta}$ and can be used to obtain bootstrap estimate of bias, standard error, variance, confidence interval. \\[0.3em]

$\bullet$ \textbf{Bootstrap estimate of Bias, standard error \& variance} of $\hat{\theta}$ are given by $$\text{Bias}_{\text{boot}} = \dfrac{1}{B} \sum \limits_{b = 1}^{B} \widehat{\theta_{b}^{*}} - \hat{\theta} \qquad \& \qquad \text{SE}_{\text{boot}} = \sqrt{\dfrac{1}{B-1} \sum \limits_{b = 1}^{B} \left[\widehat{\theta_{b}^{*}} - \overline{\theta^*}\right]^2} \qquad \& \qquad \text{var}_{\text{boot}} = \dfrac{1}{B-1} \sum \limits_{b = 1}^{B} \left[\widehat{\theta_{b}^{*}} - \overline{\theta^*}\right]^2$$

\begin{align*}
\text{where } \hat{\theta} &: \text{ estimate of } \theta \text{ from the original sample}, \\[0.2em]
\widehat{\theta_{b}^{*}} &: \text{ estimate of } \theta \text{ from } b\text{-th bootstrap sample}, \\[0.2em]
\overline{\theta^*} &: \text{ average of the bootstrap estimates}, \\[0.2em]
B &: \text{number of bootstrap samples drawn}.
\end{align*}

\section*{\faArrowAltCircleRight[regular] \textcolor{blue}{R Program}}

\begin{enumerate}

\item[\scalebox{2.0}{\twemoji{keycap: 1}}]



\begin{knitrout}
\definecolor{shadecolor}{rgb}{0.969, 0.969, 0.969}\color{fgcolor}\begin{kframe}
\begin{alltt}
\hlkwd{set.seed}\hldef{(}\hlnum{14}\hldef{)}
\end{alltt}
\end{kframe}
\end{knitrout}

\begin{knitrout}
\definecolor{shadecolor}{rgb}{0.969, 0.969, 0.969}\color{fgcolor}\begin{kframe}
\begin{alltt}
\hldef{x} \hlkwb{<-} \hlkwd{c}\hldef{(}\hlnum{2.1}\hldef{,} \hlnum{1.5}\hldef{,} \hlnum{3.2}\hldef{,} \hlnum{0.8}\hldef{,} \hlnum{1.9}\hldef{,} \hlnum{2.7}\hldef{,} \hlnum{1.1}\hldef{,} \hlnum{0.9}\hldef{)}

\hldef{theta_hat} \hlkwb{<-} \hlnum{1} \hlopt{/} \hlkwd{mean}\hldef{(x); theta_hat}
\end{alltt}
\begin{verbatim}
## [1] 0.5633803
\end{verbatim}
\end{kframe}
\end{knitrout}

\item[\scalebox{2.0}{\twemoji{keycap: 2}}]

\begin{knitrout}
\definecolor{shadecolor}{rgb}{0.969, 0.969, 0.969}\color{fgcolor}\begin{kframe}
\begin{alltt}
\hldef{n} \hlkwb{<-} \hlkwd{length}\hldef{(x); B} \hlkwb{<-} \hlnum{2000}
\end{alltt}
\end{kframe}
\end{knitrout}

\begin{knitrout}
\definecolor{shadecolor}{rgb}{0.969, 0.969, 0.969}\color{fgcolor}\begin{kframe}
\begin{alltt}
\hldef{bootstrap.theta_hat} \hlkwb{<-} \hlkwd{c}\hldef{()}

\hlkwa{for}\hldef{(i} \hlkwa{in} \hlnum{1}\hlopt{:}\hldef{B)\{}
  \hldef{bootstrap.sample} \hlkwb{<-} \hlkwd{rexp}\hldef{(n,} \hlkwc{rate} \hldef{= theta_hat)}

  \hldef{bootstrap.theta_hat[i]} \hlkwb{<-} \hlnum{1} \hlopt{/} \hlkwd{mean}\hldef{(bootstrap.sample)}
\hldef{\}}

\hlkwd{head}\hldef{(bootstrap.theta_hat)}
\end{alltt}
\begin{verbatim}
## [1] 1.2753947 1.0195017 0.7164178 0.5518297 0.3237076 0.9137179
\end{verbatim}
\end{kframe}
\end{knitrout}

\item[\scalebox{2.0}{\twemoji{keycap: 3}}]

\begin{knitrout}
\definecolor{shadecolor}{rgb}{0.969, 0.969, 0.969}\color{fgcolor}\begin{kframe}
\begin{alltt}
\hldef{bootstrap.bias} \hlkwb{<-} \hlkwd{mean}\hldef{(bootstrap.theta_hat)} \hlopt{-} \hldef{theta_hat; bootstrap.bias}
\end{alltt}
\begin{verbatim}
## [1] 0.07592151
\end{verbatim}
\end{kframe}
\end{knitrout}

\begin{knitrout}
\definecolor{shadecolor}{rgb}{0.969, 0.969, 0.969}\color{fgcolor}\begin{kframe}
\begin{alltt}
\hldef{bootstrap.sd} \hlkwb{<-} \hlkwd{sd}\hldef{(bootstrap.theta_hat); bootstrap.sd}
\end{alltt}
\begin{verbatim}
## [1] 0.2477996
\end{verbatim}
\end{kframe}
\end{knitrout}

\begin{knitrout}
\definecolor{shadecolor}{rgb}{0.969, 0.969, 0.969}\color{fgcolor}\begin{kframe}
\begin{alltt}
\hldef{bootstrap.var} \hlkwb{<-} \hlkwd{var}\hldef{(bootstrap.theta_hat); bootstrap.var}
\end{alltt}
\begin{verbatim}
## [1] 0.06140465
\end{verbatim}
\end{kframe}
\end{knitrout}


\item[\scalebox{2.0}{\twemoji{keycap: 4}}]

\begin{knitrout}
\definecolor{shadecolor}{rgb}{0.969, 0.969, 0.969}\color{fgcolor}\begin{kframe}
\begin{alltt}
\hlkwd{quantile}\hldef{(bootstrap.theta_hat,} \hlkwd{c}\hldef{(}\hlnum{0.025}\hldef{,} \hlnum{0.975}\hldef{))} \hlcom{# alpha = 0.05}
\end{alltt}
\begin{verbatim}
##      2.5%     97.5% 
## 0.3089012 1.2703311
\end{verbatim}
\end{kframe}
\end{knitrout}

\item[\scalebox{2.0}{\twemoji{keycap: 5}}]



\begin{knitrout}
\definecolor{shadecolor}{rgb}{0.969, 0.969, 0.969}\color{fgcolor}\begin{kframe}
\begin{alltt}
\hlkwd{data.frame}\hldef{(bootstrap.theta_hat)} \hlopt{%>%}
  \hlkwd{ggplot}\hldef{(}\hlkwd{aes}\hldef{(}\hlkwc{x} \hldef{= bootstrap.theta_hat))} \hlopt{+}
  \hlkwd{geom_histogram}\hldef{(}\hlkwd{aes}\hldef{(}\hlkwc{y} \hldef{=} \hlkwd{after_stat}\hldef{(density)),} \hlkwc{bins} \hldef{=} \hlnum{15}\hldef{,}
                 \hlkwc{fill} \hldef{=} \hlsng{"#275EF5"}\hldef{,} \hlkwc{color} \hldef{=} \hlsng{"black"}\hldef{,} \hlkwc{linewidth} \hldef{=} \hlnum{1}\hldef{)} \hlopt{+}
  \hlkwd{labs}\hldef{(}\hlkwc{x} \hldef{=} \hlkwd{expression}\hldef{(}\hlkwd{hat}\hldef{(theta)),}
       \hlkwc{title} \hldef{=} \hlsng{"Approximate Sampling Distribution based on Bootstrap Samples"}\hldef{)}
\end{alltt}
\end{kframe}
\includegraphics[width=\maxwidth]{figure/unnamed-chunk-11-1} 
\end{knitrout}

\smallpencil \hspace{0.1cm} {\setlength{\spaceskip}{1em plus 0.5em minus 0.5em} \fontsize{17}{20}\myfont The bootstrap samples have a positively skewed distribution. \par}

\end{enumerate}



\end{document}
